\section{Critical targets on their intrinsic residence scale}
\label{sec:v95-anchors}
A history can remain in the target core of a prescribed shell and still
have enough residence to provide its own anchor. We now use an anchor
of roof length proportional to the square of the target gradient,
on whichever physical side is longer. This does not require a crossing
of the shell prescribed in the preceding revision.

We first record the shell input with the finite-count remainder visible.
For $0<b<D<q_*$ set
\[
 v^{b,D}(u)\,du=F_\#(I_m z^{-1}\one_{\{b\le z<D\}}\bar\nu).
\]
The retained critical-volume estimate implies
\begin{equation}\label{eq:v95-shell-input}
 m^2\int_J v^{b,D}\le
 C D^{\alpha-1}(|J|+B_*^{-1})
 +C_{B_*}(1+B_*|J|)m^2b^{-1/2}\theta_{B_*}^m.
\end{equation}
Indeed multiply the positive identity for $z^{-1}\one_{b\le z<D}$
by $I_m$ and bound it by
$D^{-1}I_m\one_{z<D}+\int_b^D s^{-2}I_m\one_{z<s}\,ds$.
The two integrals are bounded by
$C D^{\alpha-1}$ and $2b^{-1/2}$, respectively. This is precisely
the finite positive layer cake of retained module 224, not an
untruncated first-inverse-gradient estimate.

Put
\[
 \eta_0=\frac1{16\Lambda_\circ},\qquad
 d_q=\frac{q^2}{256\Lambda_\circ},\qquad J_{t,q}=(t-4d_q,t+4d_q).
\]
For any original prefix $c$, let $H_{c,q}^\varepsilon$ be its actual
target source on
\[
 q\le z(t)<2q,\qquad \ell_\sigma\ge\eta_0q^2.
\]
Only targets are selected: subsequent anchor values may lie outside
the selected prefix, but remain in its same physical segment.

\begin{theorem}[Two-sided intrinsic anchor at a critical target]
\label{thm:v95-anchor}
For $3q<q_*$, every finite count and almost every target roof,
\begin{equation}\label{eq:v95-anchor-finite}
 H_{c,q}^\varepsilon(t)
 \le d_q^{-1}\int_{J_{t,q}}p_{3q}(u)\,du
       +C E_\varepsilon\int_{J_{t,q}}v^{q/2,3q}(u)\,du.
\end{equation}
In particular, at each previously fixed $q>0$,
\begin{equation}\label{eq:v95-anchor-height}
 \limsup_m\sup_{R,n,k,c}m^2\|H_{c,q}^\varepsilon\|_\infty
       \le C q^{129/128}+C E_\varepsilon q^{257/128}.
\end{equation}
There is no incidence cutoff and no required shell crossing.
One fixed-band finite-count upper bound for its normalized height is
\begin{align}
 C\{q^{\alpha-2}+E_\varepsilon q^{\alpha-1}\}(q^2+B_*^{-1})
 +C_{B_*}(1+B_*q^2)m^2\theta_{B_*}^m
                  (q^{-3/2}+E_\varepsilon q^{-1/2}).
 \label{eq:v95-anchor-remainder}
\end{align}
\end{theorem}
\begin{proof}
The selected segment has length at least $16d_q$. One side of the
target therefore has length at least $8d_q$. Choose the interval
$(t-3d_q,t-2d_q)$ on the backward side if that side is long enough;
otherwise choose $(t+2d_q,t+3d_q)$. This is a Borel choice on the
retained segment space. Each joining interval lies in $J_{t,q}$
and in the same physical segment. Radial growth gives throughout
\[
 z(u)^2\ge q^2-8\Lambda_\circ d_q>q^2/4,
 \qquad z(u)^2\le4q^2+8\Lambda_\circ d_q<9q^2.
\]
Average the elementary inequality
$f_\sigma^\varepsilon(t)\le f_\sigma^\varepsilon(v)
+\int_{[v,t]}|(f_\sigma^\varepsilon)'|$
over the anchor. Here $[v,t]$ means the unoriented joining interval.
Use $f^\varepsilon\le A$ on the first term and
$|(f^\varepsilon)'|\le C E_\varepsilon A I_m/z$ on the second.
At each intermediate roof, all joining intervals form a subfamily of
the actual physical fiber. The Borel bijection gives multiplicity one,
so their complete sum is bounded by the two original source integrals
in \eqref{eq:v95-anchor-finite}. The anchor is not a new normalized
source and does not change the target event.

Apply \eqref{eq:v94-critical-finite} and
\eqref{eq:v95-shell-input}, using $|J_{t,q}|=8d_q$.
Since $d_q$ is a fixed physical multiple of $q^2$, this gives
\eqref{eq:v95-anchor-remainder}. Fix $q,\varepsilon,B_*$ before
taking the collision limsup; only afterwards enlarge $B_*$. The
remaining terms are $Cq^\alpha+C E_\varepsilon q^{\alpha+1}$.
This proves \eqref{eq:v95-anchor-height}.
\end{proof}

There is a complementary geometric distinction which is useful at
residence zero. If
\begin{equation}\label{eq:v95-thin-condition}
                  \ell_\sigma<\eta_0 z(t)^2,
\end{equation}
then
\begin{equation}\label{eq:v95-thin-inlet}
 (z_\sigma^-)^2\ge(1-2\Lambda_\circ\eta_0)z(t)^2
                       =\tfrac78z(t)^2.
\end{equation}
Such a segment cannot have a critical backward endpoint. Conversely,
every segment born at a critical endpoint satisfies
$\ell_\sigma>t-a_\sigma\ge z(t)^2/(2\Lambda_\circ)$,
and hence belongs to the relatively long class rather than
\eqref{eq:v95-thin-condition}.

\begin{proposition}[All thin grazing-ended targets in one radial annulus]
\label{prop:v95-thin-grazing}
Let $G_{c,q}^\varepsilon$ be the original target source satisfying
$q\le z(t)<2q$, \eqref{eq:v95-thin-condition}, and having at least
one grazing physical endpoint. Then
\begin{equation}\label{eq:v95-thin-grazing}
 \limsup_m\sup_{R,n,k,c}m^2\|G_{c,q}^\varepsilon\|_\infty
                   \le C E_\varepsilon q^{257/128}.
\end{equation}
\end{proposition}
\begin{proof}
A selected segment has length at most $4\eta_0q^2$ and throughout
it $q/2<z<3q$, by the same radial inequalities. At least one
physical endpoint has zero weighted trace. The fundamental theorem
therefore bounds its target value by the absolute derivative integral
over that segment. Sum through the Borel bijection and apply
\eqref{eq:v95-shell-input} on
$(t-4\eta_0q^2,t+4\eta_0q^2)$. After the fixed-$q$ collision
limit and auxiliary-band enlargement the bound is
$C E_\varepsilon q^2q^{\alpha-1}$, as asserted.
\end{proof}

For fixed $0<a<q$ the annuli $(2^{-j-1}q,2^{-j}q]$ meeting
$[a,q)$ are a finite family. Applying the last two results and
summing their geometric leading bounds proves
\begin{equation}\label{eq:v95-core-positive-cutoff}
 \mathcal H_M\!\left(
 R_c^\varepsilon\one_{\{a\le z<q\}}
 \one_{\{\ell\ge\eta_0z^2\}\,\cup\,\{\text{a grazing endpoint}\}}
                      \right)
 \le Cq^\alpha+C E_\varepsilon q^{\alpha+1}.
\end{equation}
The notation means restriction on the physical target before
projection. The constant is independent of the previously fixed
positive $a$. The physical annulus sum remains finite before the
collision limsup. Formula \eqref{eq:v95-anchor-remainder} explains
why sending $a$ to zero inside that limit is not licensed: its
fixed-band floor and its finite-count remainder are not summable
there. The unresolved part can still concentrate in those disappearing
annuli, or on thin segments with two nongrazing endpoints.
